\section{PM-Reg Modeling}
We now discuss how to prove for a pattern $r_i$ in PM-Reg,
how does one argue that a committed string $s$
does not contain the patter $r_i$.

Our starting point will be again an AC-DFA. We collect
all patterns from a PM-Reg pattern.
Our first gadget would be to provide to Perdersen vector
commitments $\Commit_{\myvec{s}}$, $\Commit_{\myvec{i}}$,
which encodes the final state and their index along the acceptance path.

We then need to encode the patten constraints over length. One can see
that a PM-Reg is a one level boolean combination of sub-signatures, and
each sub-signautre is a sequence of fixed words and
patterns like $\ttt{.*}$ and $\ttt{.\{5,6\}}$. Given
a $(state, pos)$ sequence, one can assert that if it
matches a sub-signature by checking arithmetic constraints.
We illustrate the concept using an example.

\begin{example} Let $\ttt{abc}$ and $\ttt{def}$ and $\ttt{gh}$
be three patterns and their corresponding states be 
$s_1, s_2$, and $s_3$. Let a subsignature
$r := \ttt{abc.\{3,5\}def.\{4,6\}gh}$. Then given
a string $s := \ttt{ccabc1234def123gh}$. Running
through the ACDFA which got state pairs
$(s_1, 4), (s_2, 9), (s_3, 14)$,  one can see that
$s$ does not satisfy the regular sub-signature, because
it does not satisfy the constraints that in between
$s_2$ and $s_3$ there are $7$ to $9$ characters.

Formally, the algorithm works as the follows. Given a subsignature
it extracts the corresponding state sequence, and after
each state, there is a range, e.g., $[3,5]$ which determines the
position of the next matching state. At each step, the algorithm
receives a set of positions for the corresponding state (e.g.,
initially $\{5\}$ for state $s_1$), then it leads to allows
ranges $[6,8]$ for the appearance of second state $s_2$ in the sequence.
The algorithm then looks through the state position pairs and finds
the allowed ones. Thus, for the 2nd step, it is $(s_9, 9)$.
Then for the third state $s_3$, the allowed range is
$[15, 17]$. Then, the $(s_3, 14)$ does not fall into the range,
thus proving that $s \in L(r)$.
\end{example}

So now the problem boils down to: given a sequence of 
indices, e.g., $[1, 3, 5]$, we want to generate its subset
will falls into a set of ranges e.g., $[2,4], [5,7]$
(which generates $[3,5]$). One could actually mark up
the sorted array of these two lists appropriately:

\begin{eqnarray*}
  &r_1 & 1 2 3 4 5 5 7 \\
  &r_b & 0 1 0 0 1 0 0 \\
  &r_e & 0 0 0 1 0 0 1 \\
  &r_m & 0 1 1 1 1 1 1 \\
  &r_s & 1 0 1 0 0 1 0 \\
  &r_o & 0 0 1 0 0 1 0 
\end{eqnarray*}
Here $r_1$ represents the sorted merged list of numbers,
$r_b$ (boolean) indicates the the beginning dices,
and $r_e$ stands for the ending indices of ranges.
Note that $r_b$ and $r_e$ can be checked by 
finger printing techniques (if it's consistent with
a subsignature pattern).
Then $r_m$ marks all elements that are in some ranges,
this can be accomplished by checking formula that
$r_m[i] = 1$ if $r_e[i] \neq 1$ and its previous element is $1$
or $r_b[i-1] = 1$. 

$r_s$ represents the index of elements in input sequence.
$r_o$ represents the output , which is a bit-wise
conjunction between $r_s$ and $r_s$.

Altogether the prover complexity is linear over the sum of
the number of states and patterns. Given that the average of
final states is $6.6k$, the concrete cost is $10\times$, i.e.,
$66k$. Given max 50 patterns, this is $3$ million constraints max
and for average 4 patterns, this is only $240k$.

\subsection{Boolean Combination of Subsignatures}
A regex pattern is a boolean combination of sub-signatures, in the
CNF form. Then each conjunction can be represented by a bit pattern,
e.g., assuming 5 subsignatures in max, it consists of 10 bits.
These bit-patterns (converted to field elements) are stored in a lookup
table.
At prover time, the prover feeds these bit-patterns for a PM-REG pattern
in a separate advice column, and they are verified using lookup argument.
One field element can store 10 20-bit patterns and this is sufficient
for our application.

Now given a CNF pattern $\myvec{a}$ and the corresponding
subsignature evaluation $\myvec{b}$, the evaluation could be
defined as the result $c = \prod_{i\in [10]} (\myvec{a}_i=0 \wedge 1) \vee
	(\myvec{a}_i=1 \wedge \myvec{b}_i)$.

\subsection{Question:}
{\bf Do a global optimizer or using gadgets in halo2?}



